% AY2025 制御工学 第1問〔I〕（転記）
% 出典: 03_制御工学/original/03_制御工学/2025/2025se.pdf
% 要約: 壁-並列(k2,d2)-質量 m-並列(k1,d1)-入力端 A(変位 x1)。
%       定常値・ステップ応答・最大値、(5)で A 固定し質量に外力 f(t)。

\section*{第1問〔I〕}

図1に示す系に関する次の問い (1)〜(5) に答えよ. なお, 質量を $m$, 粘性係数を $d_1$, $d_2$,
バネ定数を $k_1$, $k_2$ とし, $x_1(t)$ は系への入力となる変位, $x_2(t)$ は質量の変位を表す.
ダッシュポットについては動作速度に比例した粘性抵抗が生じるものとし, 初期状態において
系は静止しているとする.

\begin{center}
\begin{tikzpicture}[>=Latex]
  % 壁
  \draw[thick] (0,0.3) -- (0,2.0);
  \foreach \y in {0.4,0.55,...,1.95} \draw[thick] (0,\y) -- (-0.22,\y-0.22);
  % 左 並列: バネ k2（上）+ ダンパ d2（下）
  \draw[thick,decorate,decoration={zigzag,segment length=7,amplitude=4}] (0,1.55) -- (1.6,1.55);
  \node at (0.8,1.95) {$k_2$};
  \draw[thick] (0,0.85) -- (0.45,0.85);
  \draw[thick] (0.45,0.63) rectangle (0.9,1.07);
  \draw[thick] (0.7,0.85) -- (1.6,0.85);
  \node at (0.8,0.35) {$d_2$};
  \draw[thick] (1.6,1.55) -- (1.6,0.85);
  % 質量 m
  \draw[thick,fill=gray!12] (1.6,0.45) rectangle (2.35,1.95);
  \node at (1.975,1.2) {$m$};
  \draw[->,thick] (1.975,0.2) -- (2.6,0.2);
  \node at (2.7,-0.02) {$x_2(t)$};
  % 右 並列: バネ k1（上）+ ダンパ d1（下）
  \draw[thick,decorate,decoration={zigzag,segment length=7,amplitude=4}] (2.35,1.55) -- (3.95,1.55);
  \node at (3.15,1.95) {$k_1$};
  \draw[thick] (2.35,0.85) -- (2.8,0.85);
  \draw[thick] (2.8,0.63) rectangle (3.25,1.07);
  \draw[thick] (3.05,0.85) -- (3.95,0.85);
  \node at (3.15,0.35) {$d_1$};
  \draw[thick] (3.95,1.55) -- (3.95,0.85);
  % 入力端 A
  \draw[thick] (3.95,1.2) -- (4.25,1.2);
  \draw[thick,fill=white] (4.25,1.2) circle (0.06);
  \node at (4.55,1.2) {A};
  \draw[->,thick] (4.25,0.95) -- (4.9,0.95);
  \node at (4.75,0.72) {$x_1(t)$};
  \node at (2.2,-0.45) {\small 図1};
\end{tikzpicture}
\end{center}

\begin{enumerate}[label=(\arabic*)]
  \item $x_1(t)$ にステップ入力を与えたときの $x_2(t)$ の定常値を求めよ.

  \item 各パラメータを $m=1$, $d_1=1$, $d_2=1$, $k_1=1$, $k_2=2$ とし, 入力として $x_1(t)$ に
        ステップ入力を与えたときの応答 $x_2(t)$ を求めよ. また, 定常値が存在する場合には,
        定常値についても求めよ.

  \item 問い (2) で求めた応答 $x_2(t)$ の最大値とその発生時間を求めよ.

  \item 各パラメータを $m=1$, $d_1=3$, $d_2=0$, $k_1=1$, $k_2=1$, 初期値をすべて $0$ とし,
        入力として $x_2(t)$ にステップ入力を与えたときの応答 $x_2(t)$ と最大値の発生時間を
        求めよ.

  \begin{rem}
  原本では本小問の「入力として $x_2(t)$ に……」と記載されているが, 系の入力は $x_1(t)$ で
  あり, $x_2(t)$ は出力であるため, 入力は $x_1(t)$ の誤記と思われる.
  \end{rem}

  \item 次に, すべてのバネが自然長の状態において点 A を固定端に取付け, 各パラメータを
        $m=1$, $d_1=1$, $d_2=1$, $k_1=1$, $k_2=1$ とし, 質量にステップ入力となる力 $f(t)$ を
        加える. この条件下において, 初期値をすべて $0$ としたときの応答 $x_2(t)$ を求めよ.
\end{enumerate}
